%% ma: L11-B09-0005
%% lop: 11
%% chuong: 3
%% bai: 9
%% dang: 11.9.4
%% loai: TN4
%% muc: TH
%% dap_an: "C"
%% nguon: "chuyen-de-cac-so-dac-trung-cua-mau-so-lieu-khong-ghep-nhom-toan-10-chuong-trinh-moi.pdf (GV. Phan Nhật Linh) (P02), Bài 02 Dạng 1 Phần I Câu 35"
%% kien_thuc: "Lớp 11 Bài 9 (số trung bình của mẫu số liệu ghép nhóm)"
%% trang_thai: chinh-thuc
%% ghi_chu: "Mẫu số liệu ghép nhóm (vượt Lớp 10 Bài 13), đã chuyển sang Lớp 11 Bài 9. Tự duyệt theo yêu cầu giáo viên 22/10/2026."
\begin{ex}
Một người lái xe thường xuyên đi lại giữa hai thành phố $A$ và $B$. Thời gian đi (phút) được ghi lại trong bảng tần số ghép nhóm sau:
\begin{center}
\begin{tabular}{|c|c|c|c|c|c|c|c|}\hline
Thời gian & $[40;44]$ & $[45;49]$ & $[50;54]$ & $[55;59]$ & $[60;64]$ & $[65;69]$ & Cộng\\\hline
Tần số & $9$ & $15$ & $30$ & $17$ & $17$ & $12$ & $100$\\\hline
\end{tabular}
\end{center}
Thời gian trung bình người đó đi từ thành phố $A$ đến $B$ là
\choice{$54{,}5$ phút}{$54{,}6$ phút}{\True $54{,}7$ phút}{$54{,}8$ phút}
\loigiai{Giá trị đại diện: $42$; $47$; $52$; $57$; $62$; $67$.

$\overline x=\dfrac{42\cdot9+47\cdot15+52\cdot30+57\cdot17+62\cdot17+67\cdot12}{100}=\dfrac{5470}{100}=54{,}7$. Chọn C.}
\end{ex}
